\documentclass{article}
\usepackage{amsmath,amsthm,amssymb}
\usepackage{algorithm}
\usepackage{algpseudocode}
\usepackage{hyperref}
\usepackage{enumitem}
\newtheorem{theorem}{Theorem}
\newtheorem{lemma}{Lemma}
\newtheorem{assumption}{Assumption}
\newcommand{\argmin}{\operatorname*{arg\,min}}

\begin{document}

%%%%%%%%%%%%%%%%%%%%%%%%%%%%%%%%%%%%%%%%%%%%%%%%%%%%%%%%%%%%%%%%%%%%%%%%%%%%%%
\section*{Algorithm Name}
\textbf{Frank–Wolfe (Conditional Gradient) Method}
%%%%%%%%%%%%%%%%%%%%%%%%%%%%%%%%%%%%%%%%%%%%%%%%%%%%%%%%%%%%%%%%%%%%%%%%%%%%%%

%%%%%%%%%%%%%%%%%%%%%%%%%%%%%%%%%%%%%%%%%%%%%%%%%%%%%%%%%%%%%%%%%%%%%%%%%%%%%%
\section*{Mathematical Formulation}
Let $f:\mathbb{R}^n\rightarrow\mathbb{R}$ be a convex and $L$‑smooth function and let 
$C\subset\mathbb{R}^n$ be a non‑empty compact convex feasible set.  
Given the current iterate $x_k\in C$ and a stepsize $\gamma_k\in[0,1]$, the method performs

\[
\begin{aligned}
s_k &:= \argmin_{s\in C}\;\big\langle \nabla f(x_k),\,s\big\rangle ,\tag{1}\\[4pt]
x_{k+1} &:= (1-\gamma_k)\,x_k+\gamma_k\,s_k .\tag{2}
\end{aligned}
\]

The linear minimisation oracle (LMO) in (1) returns the feasible point that minimizes the
linearisation of $f$ at $x_k$.  The stepsize may be chosen by an exact line search,
by the classic FW schedule $\gamma_k=\frac{2}{k+2}$, or by any rule satisfying
$0<\gamma_k\le 1$.

%%%%%%%%%%%%%%%%%%%%%%%%%%%%%%%%%%%%%%%%%%%%%%%%%%%%%%%%%%%%%%%%%%%%%%%%%%%%%%
\section*{Pseudocode}
\begin{algorithm}[H]
\caption{Frank–Wolfe (Conditional Gradient)}
\begin{algorithmic}[1]
\State \textbf{Input:} convex $f$, compact convex $C$, initial $x_0\in C$, stepsize rule $\{\gamma_k\}$, max iter $K$.
\For{$k=0,1,\dots,K-1$}
    \State Compute gradient $g_k \gets \nabla f(x_k)$.
    \State \textbf{LMO:} $s_k \gets \argmin_{s\in C}\langle g_k , s\rangle$.
    \State Choose stepsize $\gamma_k\in[0,1]$ (e.g. $\gamma_k=2/(k+2)$ or line search).
    \State Update $x_{k+1} \gets (1-\gamma_k)x_k + \gamma_k s_k$.
    \State \textbf{Optional:} compute dual gap $g_k^{\text{FW}}:=\langle g_k, x_k-s_k\rangle$ for stopping.
\EndFor
\State \textbf{Output:} $x_K$ (or the best iterate according to the dual gap).
\end{algorithmic}
\end{algorithm}
%%%%%%%%%%%%%%%%%%%%%%%%%%%%%%%%%%%%%%%%%%%%%%%%%%%%%%%%%%%%%%%%%%%%%%%%%%%%%%

%%%%%%%%%%%%%%%%%%%%%%%%%%%%%%%%%%%%%%%%%%%%%%%%%%%%%%%%%%%%%%%%%%%%%%%%%%%%%%
\section*{Dry Run Example}
We minimise the one‑dimensional quadratic 
\[
f(w)=(w-3)^2,
\]
over the interval $C=[0,5]$.
The gradient is $\nabla f(w)=2(w-3)$.  
We initialise $w_0=0$ and use the classic FW stepsize $\gamma_k=\frac{2}{k+2}$.
The LMO in one dimension is trivial:
\[
s_k=
\begin{cases}
0, & \text{if } \nabla f(w_k)\ge 0,\\[2pt]
5, & \text{if } \nabla f(w_k)<0 .
\end{cases}
\]

\begin{center}
\begin{tabular}{c|c|c|c|c}
$k$ & $w_k$ & $\nabla f(w_k)$ & $s_k$ & $w_{k+1}= (1-\gamma_k)w_k+\gamma_k s_k$\\ \hline
0 & $0$ & $-6$ & $5$ & $\displaystyle w_1=(1-\tfrac{2}{2})\cdot0+\tfrac{2}{2}\cdot5=5$ \\[6pt]
1 & $5$ & $+4$ & $0$ & $\displaystyle w_2=(1-\tfrac{2}{3})\cdot5+\tfrac{2}{3}\cdot0=\frac{5}{3}\approx1.667$ \\[6pt]
2 & $1.667$ & $-2.667$ & $5$ & $\displaystyle w_3=(1-\tfrac{2}{4})\cdot1.667+\tfrac{2}{4}\cdot5= \frac{1}{2}\cdot1.667+ \frac{1}{2}\cdot5\approx3.333$
\end{tabular}
\end{center}

Objective values:

\[
f(w_0)=9,\quad f(w_1)=4,\quad f(w_2)=\bigl(\tfrac{5}{3}-3\bigr)^2\approx1.78,\quad 
f(w_3)=\bigl(3.333-3\bigr)^2\approx0.11 .
\]

The iterates quickly approach the minimiser $w^\star=3$.

%%%%%%%%%%%%%%%%%%%%%%%%%%%%%%%%%%%%%%%%%%%%%%%%%%%%%%%%%%%%%%%%%%%%%%%%%%%%%%
\section*{Assumptions}
\begin{assumption}[Convexity and Smoothness]
\label{as:convex-smooth}
The objective $f$ is convex and $L$‑smooth on $\mathbb{R}^n$, i.e.
\[
\|\nabla f(x)-\nabla f(y)\| \le L\|x-y\|,\qquad\forall x,y\in\mathbb{R}^n .
\]
\end{assumption}

\begin{assumption}[Compact Feasible Set]
\label{as:compact}
$C\subset\mathbb{R}^n$ is non‑empty, convex, compact and has finite diameter
\[
D:=\max_{x,y\in C}\|x-y\|<\infty .
\]
\end{assumption}

\begin{assumption}[Stepsize Rule]
\label{as:stepsize}
The stepsizes satisfy $\gamma_k\in[0,1]$ and either
\begin{enumerate}[label=(\alph*)]
\item \textbf{Exact line search:}
\[
\gamma_k = \argmin_{\gamma\in[0,1]} f\big((1-\gamma)x_k+\gamma s_k\big);
\]
\item \textbf{FW schedule:} $\displaystyle\gamma_k=\frac{2}{k+2}\,.$
\end{enumerate}
\end{assumption}

%%%%%%%%%%%%%%%%%%%%%%%%%%%%%%%%%%%%%%%%%%%%%%%%%%%%%%%%%%%%%%%%%%%%%%%%%%%%%%
\section*{Main Convergence Theorem}
\begin{theorem}[Primal gap convergence]
\label{thm:fw-rate}
Under Assumptions \ref{as:convex-smooth}–\ref{as:stepsize},
let $x^\star\in\arg\min_{x\in C}f(x)$ and define the primal gap $h_k:=f(x_k)-f(x^\star)$.  
Then for all $k\ge0$,
\[
h_k \le \frac{2L D^{2}}{k+2}.
\tag{3}
\]
Consequently $\displaystyle \lim_{k\to\infty} h_k =0$ and the rate is $O(1/k)$.
\end{theorem}

A more refined bound uses the curvature constant $C_f$ (see Lemma \ref{lem:curvature}) and yields
\[
h_k \le \frac{2C_f}{k+2}.
\]

%%%%%%%%%%%%%%%%%%%%%%%%%%%%%%%%%%%%%%%%%%%%%%%%%%%%%%%%%%%%%%%%%%%%%%%%%%%%%%
\section*{Theoretical Properties}
\begin{itemize}
\item \textbf{Stability:} The iterates always remain in $C$ because the update is a convex combination of two points in $C$.
\item \textbf{Hyperparameter Sensitivity:} The algorithm requires only a stepsize rule; the classic $\gamma_k=2/(k+2)$ guarantees the $O(1/k)$ bound without any tuning.
\item \textbf{Comparison to Gradient Descent:} When projections onto $C$ are expensive, FW avoids them at the cost of a slower $O(1/k)$ rate (vs. $O(1/k)$ for projected GD with the same smoothness but with a cheaper per‑iteration cost).
\item \textbf{Special Cases:}
    \begin{enumerate}
        \item If $f$ is linear, the algorithm terminates in one iteration.
        \item If the optimum lies on a face of $C$, the dual gap $g_k^{\text{FW}}$ provides a certificate of optimality.
    \end{enumerate}
\end{itemize}
%%%%%%%%%%%%%%%%%%%%%%%%%%%%%%%%%%%%%%%%%%%%%%%%%%%%%%%%%%%%%%%%%%%%%%%%%%%%%%
\section*{Preliminaries (Definitions and Lemmas)}
\begin{definition}[Curvature constant]
For a convex and $L$‑smooth $f$, the curvature over $C$ is
\[
C_f := \sup_{\substack{x,s\in C\\ \gamma\in[0,1]}}
\frac{2}{\gamma^{2}}\Bigl(f\big((1-\gamma)x+\gamma s\big)-f(x)
-\gamma\langle\nabla f(x),s-x\rangle\Bigr).
\tag{4}
\]
\end{definition}

\begin{lemma}[Smoothness upper bound]
\label{lem:smooth}
For any $x,y\in\mathbb{R}^n$,
\[
f(y) \le f(x)+\langle\nabla f(x),y-x\rangle+\frac{L}{2}\|y-x\|^{2}.
\tag{5}
\]
\end{lemma}
\begin{proof}
Standard consequence of $L$‑smoothness; see e.g. Nesterov (2004). \qedhere
\end{proof}

\begin{lemma}[Curvature bound]
\label{lem:curvature}
For a $L$‑smooth function over a set of diameter $D$, the curvature satisfies
\[
C_f \le L D^{2}.
\tag{6}
\]
\end{lemma}
\begin{proof}
From (5) with $y=(1-\gamma)x+\gamma s$ and using $\|s-x\|\le D$, we obtain
\[
f\big((1-\gamma)x+\gamma s\big) \le f(x)+\gamma\langle\nabla f(x),s-x\rangle
+\frac{L\gamma^{2}}{2}D^{2}.
\]
Rearranging and multiplying by $2/\gamma^{2}$ yields (6). \qedhere
\end{proof}
%%%%%%%%%%%%%%%%%%%%%%%%%%%%%%%%%%%%%%%%%%%%%%%%%%%%%%%%%%%%%%%%%%%%%%%%%%%%%%
\section*{Main Proof (Step‑by‑Step Derivation)}
We prove Theorem \ref{thm:fw-rate}.  
All non‑trivial steps are numbered for easy reference.

\subsection*{Step 1: Descent Lemma applied to the FW update}
From Lemma \ref{lem:smooth} with $x=x_k$ and $y=x_{k+1}=(1-\gamma_k)x_k+\gamma_k s_k$,
\begin{align}
f(x_{k+1}) &\le f(x_k)+\langle\nabla f(x_k),x_{k+1}-x_k\rangle
+\frac{L}{2}\|x_{k+1}-x_k\|^{2} \nonumber\\
&= f(x_k)+\gamma_k\langle\nabla f(x_k),s_k-x_k\rangle
+\frac{L\gamma_k^{2}}{2}\|s_k-x_k\|^{2}. \tag{7}
\end{align}
\textbf{[Step 1]}

\subsection*{Step 2: Use optimality of the LMO}
By definition of $s_k$,
\[
\langle\nabla f(x_k),s_k\rangle \le \langle\nabla f(x_k),x^\star\rangle,
\]
hence
\[
\langle\nabla f(x_k),s_k-x_k\rangle
\le \langle\nabla f(x_k),x^\star-x_k\rangle
= -\bigl(f(x_k)-f(x^\star)\bigr) -\underbrace{\bigl(f(x^\star)-\langle\nabla f(x_k),x^\star\rangle\bigr)}_{\ge 0\text{ by convexity}}.
\]
Thus,
\[
\langle\nabla f(x_k),s_k-x_k\rangle \le -h_k . \tag{8}
\]
\textbf{[Step 2]}

\subsection*{Step 3: Bound the distance term}
Since $s_k,x_k\in C$, by definition of the diameter $D$,
\[
\|s_k-x_k\|\le D . \tag{9}
\]
\textbf{[Step 3]}

\subsection*{Step 4: Combine (7)–(9)}
Insert (8) and (9) into (7):
\[
f(x_{k+1}) \le f(x_k) -\gamma_k h_k + \frac{L D^{2}}{2}\gamma_k^{2}. \tag{10}
\]
Subtract $f(x^\star)$ from both sides and denote $h_{k+1}=f(x_{k+1})-f(x^\star)$:
\[
h_{k+1} \le (1-\gamma_k)h_k + \frac{L D^{2}}{2}\gamma_k^{2}. \tag{11}
\]
\textbf{[Step 4]}

\subsection*{Step 5: Choose the classic FW stepsize}
Set $\displaystyle\gamma_k=\frac{2}{k+2}$ (which satisfies $0<\gamma_k\le1$ for all $k\ge0$).  
Plugging into (11) gives
\[
h_{k+1} \le \Bigl(1-\frac{2}{k+2}\Bigr)h_k + \frac{L D^{2}}{2}\Bigl(\frac{2}{k+2}\Bigr)^{2}
= \frac{k}{k+2}h_k + \frac{2L D^{2}}{(k+2)^{2}}. \tag{12}
\]
\textbf{[Step 5]}

\subsection*{Step 6: Induction on the bound}
We prove by induction that
\[
h_k \le \frac{2L D^{2}}{k+2},\qquad\forall k\ge0. \tag{13}
\]

\paragraph{Base case $k=0$:}  
From convexity, $h_0 = f(x_0)-f(x^\star) \le \frac{L}{2}D^{2}\le \frac{2L D^{2}}{2}$, so (13) holds.

\paragraph{Inductive step:}  
Assume (13) holds for some $k$. Using (12),
\[
h_{k+1}\le \frac{k}{k+2}\cdot\frac{2L D^{2}}{k+2}
      +\frac{2L D^{2}}{(k+2)^{2}}
      =\frac{2L D^{2}}{(k+2)^{2}}\Bigl(k+1\Bigr)
      =\frac{2L D^{2}}{(k+1)+2}.
\]
Thus (13) holds for $k+1$. By induction, (13) is true for all $k$.

\textbf{[Step 6]}

\subsection*{Step 7: Conclude the rate}
Equation (13) is exactly the statement of Theorem \ref{thm:fw-rate}:
\[
h_k \le \frac{2L D^{2}}{k+2}=O\!\left(\frac{1}{k}\right).
\]
\textbf{[Step 7]}

\subsection*{Step 8: Extension via curvature}
If we replace $L D^{2}$ by the curvature $C_f$ (Lemma \ref{lem:curvature}), the same induction yields
\[
h_k \le \frac{2C_f}{k+2}.
\tag{14}
\]
\textbf{[Step 8]}

%%%%%%%%%%%%%%%%%%%%%%%%%%%%%%%%%%%%%%%%%%%%%%%%%%%%%%%%%%%%%%%%%%%%%%%%%%%%%%
\section*{Rate Derivation (With Constants)}
Collecting the constants from the proof:

\[
\boxed{\;h_k \le \underbrace{\frac{2C_f}{k+2}}_{\text{primal gap bound}}\;}
\qquad\text{with } C_f \le L D^{2}.
\]

If an exact line search is used, the same recursion (11) holds with the optimal
$\gamma_k$, and the bound improves to $h_k\le \frac{2C_f}{k+2}$ without the factor $2$ in the numerator (see Jaggi, 2013).

%%%%%%%%%%%%%%%%%%%%%%%%%%%%%%%%%%%%%%%%%%%%%%%%%%%%%%%%%%%%%%%%%%%%%%%%%%%%%%
\section*{Special Cases}
\begin{enumerate}[label=\arabic*.]
\item \textbf{Polytope with vertices as extreme points.}  
When $C$ is a polytope, each $s_k$ is a vertex, and the algorithm can be interpreted as a sparse representation of $x_k$ in terms of vertices.
\item \textbf{Strongly convex $f$.}  
If $f$ is $\mu$‑strongly convex, one can combine the $O(1/k)$ primal bound with a linear rate on the dual gap, yielding $h_k = O\bigl((1-\frac{\mu}{L})^{k}\bigr)$ after a finite number of “away steps’’ (the Frank–Wolfe with away steps algorithm).
\item \textbf{Exact solution of the LMO.}  
If the LMO can be solved analytically (as in the dry‑run example), the per‑iteration cost is negligible, making FW attractive for large‑scale constrained problems.
\end{enumerate}

%%%%%%%%%%%%%%%%%%%%%%%%%%%%%%%%%%%%%%%%%%%%%%%%%%%%%%%%%%%%%%%%%%%%%%%%%%%%%%
\section*{Discussion}
The Frank–Wolfe method enjoys a simple implementation (only a linear oracle) and provable $O(1/k)$ convergence under the mild assumptions of convexity, smoothness, and compactness.  
The proof above follows directly from the smoothness inequality, the optimality of the LMO, and a careful choice of stepsize.  
All constants are explicit, allowing practitioners to estimate the number of iterations required to achieve a desired accuracy.

\end{document}